\section{Introduction}
\label{sec:introduction}

An LLM serving request cannot generate a token until model weights are resident. In serverless and multi-tenant settings, idle models are commonly evicted to recover scarce accelerator memory; the next request then pays a cold load before inference begins. A 7B float16 checkpoint is already on the order of 14~GB, and the checkpoints in this paper reach beyond 65~GB of weight bytes. Cold loading is therefore not a startup footnote; it is part of the user-visible serving path.

The usual shorthand is to call this a bandwidth problem. That shorthand is incomplete. A loader does not read an abstract number of gigabytes; it executes a concrete read plan. The checkpoint format decides whether bytes appear as a few contiguous shards or as many indexed ranges across partition files. The I/O backend decides whether those reads are submitted through blocking POSIX calls, an async runtime, or \texttt{io\_uring}. The integration layer then adds its own costs: Python object construction, tensor materialization, CUDA context setup, allocator initialization, and serving-engine configuration. A device can have high peak throughput and still produce poor cold-start time when request geometry and submission method are mismatched.

Existing LLM serving systems and checkpoint libraries typically treat loading as a fixed initialization cost with a single code path. No prior work, to our knowledge, systematically evaluates I/O backend choice for LLM checkpoint loading as a function of format, scale, and deployment environment. This gap matters because the optimal backend is not stable across the parameter space.

This paper evaluates three explicit I/O backends: synchronous POSIX I/O, Tokio-based async loading, and Linux \texttt{io\_uring}, across two checkpoint layouts, six model scales, four environments, and two measured layers (Rust and vLLM) with Python as the integration bridge. Our central finding is:

\begin{quote}
\textbf{No single I/O backend dominates.} The best backend depends on checkpoint layout, model scale, backend availability, and the layer at which loading is observed. A fixed rule such as ``always sync,'' ``always async,'' or ``always \texttt{io\_uring}'' leaves performance on the table.
\end{quote}

SafeTensors presents mostly whole-file or multi-shard eager loading: after metadata parsing, the loader must move large contiguous byte regions. ServerlessLLM-style partitioning presents a different workload: tensors are reconstructed from indexed partition files, creating many grouped range reads. These layouts are not cosmetic alternatives; they expose different I/O regimes. Scale then moves the boundary within a regime: a backend with low orchestration overhead can win for a small single-shard load, while a backend that sustains deeper submission queues can win for a large multi-shard load.

We built Tensora\footnote{Tensora is this paper's checkpoint-loading framework (Rust crate \texttt{tensora}, Python package \texttt{tensora\_py}). It is unrelated to Google's TensorStore library for array storage.}, an open framework for evaluating these regimes. Tensora exposes sync, async, and \texttt{io\_uring} under shared format abstractions, enabling controlled backend attribution across layers. We evaluate six public checkpoints from 0.6B to 32B parameters across two layouts and two measured layers: Rust backend attribution and a custom vLLM integration. The primary environment is an H100 host with local NVMe and working \texttt{io\_uring}; H200 and A100 container environments with network-attached storage and blocked \texttt{io\_uring} test whether ordering regimes transfer when the strongest H100 backend is unavailable. A Modal H100 environment under gVisor provides restricted-container validation.

On H100 SafeTensors, sync wins at the small end and \texttt{io\_uring} from Qwen3-8B through Qwen3-32B (1.28$\times$ over sync). On H100 ServerlessLLM, sync is never fastest; async and \texttt{io\_uring} lead different rows. vLLM reorders backend rankings (it measures a different construct) but the effect survives integration: in vLLM cold \texttt{load\_only}, Tensora's \texttt{io\_uring} SafeTensors path reduces initialization by 37 to 48\% relative to vLLM's native loader for 8B and 32B checkpoints. Cross-environment results preserve small-scale sync dominance on three of four environments while showing that the large-scale winner depends on kernel and storage context.

\textbf{Contributions.}
\begin{enumerate}
  \item \textbf{Format-driven backend dependence.} We evaluate three I/O backends (sync, async, \texttt{io\_uring}) across two checkpoint layouts and six model scales, showing that the best backend shifts systematically with format and scale: sync wins on small SafeTensors, \texttt{io\_uring} wins on large SafeTensors, and sync never wins on H100 ServerlessLLM. To our knowledge, this is the first controlled evidence that LLM checkpoint loading is workload-dependent rather than bandwidth-bound.
  \item \textbf{Backend regime persistence through integration.} We provide controlled evidence across two measured layers (Rust and vLLM) demonstrating that backend ordering regimes propagate through library and serving-engine boundaries. The proportional advantage shrinks but the regime structure survives.
  \item \textbf{Regime transfer across environments.} We validate across four environments (H100, H200, A100, Modal) that small-scale sync dominance transfers widely while the large-scale winner depends on storage topology and kernel capability. Regime structure, not backend identity, is the portable quantity.
\end{enumerate}
